\section{ArgoCD on the OpenShift Developer Sandbox}

The original assignment included setting up ArgoCD to handle GitOps-style deployments,
where every commit to the infrastructure repository would automatically update what is
running in the cluster. In the original setup with the shared cluster this would have
been straightforward. The Sandbox is a different situation.

\subsection{What ArgoCD Needs}

ArgoCD is not a simple application. When it is installed in a cluster it registers
several cluster-wide Custom Resource Definitions, CRDs, that define new resource types
like \texttt{Application} and \texttt{AppProject}. It also creates service accounts with
permissions to read and write resources across the entire cluster, not just one namespace.
The \texttt{argocd cluster add} CLI command, which registers an external cluster for
ArgoCD to manage, works by creating a service account in the \texttt{kube-system}
namespace of the target cluster.

The Sandbox does not allow any of this. Users have access only to their own namespace.
CRD creation is blocked. RBAC escalation is blocked. The \texttt{kube-system} namespace
is inaccessible entirely. These are not bugs in the Sandbox but deliberate restrictions
to prevent one tenant from affecting another.

\subsection{What I Tried}

I went through twelve distinct approaches over several sessions trying to get ArgoCD
working, or at least partially working. None of them fully succeeded inside the Sandbox,
but each one taught me something concrete about how ArgoCD and OpenShift work.

The first attempt was to install the OpenShift GitOps Operator through OperatorHub,
which is the official way to get ArgoCD on an OpenShift cluster. The OperatorHub menu
is simply not available in the Sandbox, so that path was closed immediately.

The second attempt was to apply the official \texttt{namespace-install.yaml} manifest
directly. This failed because the manifest tries to create CRDs and escalate RBAC
privileges, both of which are blocked.

For the third attempt I set up a local Kind cluster on my Mac using Podman as the
container runtime. This worked after I set the environment variable
\texttt{KIND\_EXPERIMENTAL\_PROVIDER=podman}, which is needed because Kind defaults to
Docker. I then installed ArgoCD v3.3.2 into that Kind cluster successfully. Having a
local ArgoCD instance to work with was useful for the rest of the attempts even though
it was not the target environment.

Attempts five and six tried to register the Sandbox cluster as a target in my local
ArgoCD using \texttt{argocd cluster add} with first a cluster-level and then a
namespace-scoped flag. Both failed for the same reason: the CLI tries to create a
service account in \texttt{kube-system} on the target cluster regardless of the flags.

Attempts seven and eight tried to work around this by manually creating the
\texttt{argocd-manager} service account in my own namespace. The first try used admin
privileges, which RBAC blocked. The second try used edit privileges, which the cluster
accepted, but edit access to a single namespace is not enough for ArgoCD to perform a
cluster sync.

Attempt nine created a manual cluster secret in ArgoCD containing my own user token.
ArgoCD accepted it and the Sandbox cluster showed up in the interface. However when I
clicked Sync it failed because my user token does not have cluster-scope read
permissions, only namespace-scope.

Attempt ten tried to use a service account token instead of the user token. Service
account tokens in OpenShift 4.11 and later are not long-lived by default and have to be
requested differently than in older versions. The cluster-scope access problem was still
present regardless.

Attempt eleven applied ArgoCD's own manifests directly inside my Sandbox namespace.
Several pods came up but others failed due to Security Context Constraints blocking the
privilege levels that ArgoCD components require.

Attempt twelve tried cross-cluster token sharing between the Kind cluster and the
Sandbox. Tokens are signed by the cluster that issued them and are not valid on other
clusters, so this was a dead end.

\subsection{What I Did Achieve}

The application stack itself was fully deployed and working throughout all of this. The
frontend was reachable over HTTPS through the Route, the backend was connecting to both
PostgreSQL and Redis, and persistent storage was functioning correctly for PostgreSQL.
ArgoCD v3.3.2 was installed and running in my local Kind cluster. The Sandbox cluster
was successfully registered in that ArgoCD instance using a manual cluster secret, which
is actually a technique used in production multi-cluster setups. The only thing missing
was the synchronization step, which requires permissions the Sandbox will not grant.

\subsection{What This Exercise Actually Taught Me}

Going through twelve failing attempts was more educational than a clean successful
install would have been. I now understand exactly which permissions ArgoCD needs and
why: it has to enumerate resources across the whole cluster to know what state things
are in, and it has to write to multiple namespaces to apply manifests. A namespace-scoped
installation simply cannot do the same job.

For anyone who needs GitOps on the Sandbox the practical answer is to use GitHub Actions
instead. The CI workflow can apply manifests directly with \texttt{oc apply} using a
token scoped to the project namespace, which is something the Sandbox does permit. It is
not as elegant as ArgoCD but it works within the constraints of the environment.

If I were doing this on a real cluster, either the shared \texttt{devops24} cluster or
a managed cloud cluster on EKS, AKS, or GKE, all of these attempts would have
succeeded. The exercise showed me clearly where the boundaries are.
